%% pc-echelle-ph — échelle de pH de 0 à 14 et concentration en ions oxonium.
%% La bande colorée est une suite de rectangles d'une unité de pH (rouge -> vert -> bleu,
%% couleurs d'un indicateur universel) : pas de `shading`, le SVG reste léger.
%% Règle du bas : [H3O+] = 10^(-pH), donc la concentration DIMINUE quand le pH augmente.
\documentclass[tikz,border=4pt]{standalone}
\usepackage{liaisons}
\begin{document}
\begin{tikzpicture}[x=0.5cm, y=1cm,
  axe/.style={-{Stealth[length=5pt]}, line width=0.6pt},
  grad/.style={font=\scriptsize, inner sep=2pt},
  petit/.style={font=\scriptsize, inner sep=1.5pt},
  trait/.style={line width=0.4pt, black!55}]

\def\hb{0.55}   % hauteur de la bande colorée

%% ---------------- bande colorée : 14 rectangles d'une unité de pH ----------------
\foreach \k in {0,1,...,6} {                     % acides : rouge -> vert
  \pgfmathsetmacro\p{100-100*(\k+0.5)/7}
  \fill[solideA!\p!solideE] (\k,0) rectangle (\k+1,\hb); }
\foreach \k in {7,8,...,13} {                    % basiques : vert -> bleu
  \pgfmathsetmacro\p{100-100*(\k+0.5-7)/7}
  \fill[solideE!\p!solideB] (\k,0) rectangle (\k+1,\hb); }
\draw[line width=0.6pt] (0,0) rectangle (14,\hb);
\draw[line width=0.9pt, white] (7,0) -- (7,\hb);

%% ---------------- domaines acide / basique ---------------------------------------
\node[font=\small\bfseries, text=solideA] at (3.5,1.62) {solutions acides};
\node[font=\small\bfseries, text=solideB] at (10.5,1.62) {solutions basiques};
\draw[axe, solideA] (6.6,1.30) -- (0.2,1.30);
\draw[axe, solideB] (7.4,1.30) -- (13.8,1.30);

%% ---------------- trois repères usuels -------------------------------------------
\foreach \p/\lab in {2/{jus de citron}, 7/{eau pure}, 13/{déboucheur}} {
  \draw[trait] (\p,\hb) -- (\p,0.98);
  \node[petit, anchor=south, align=center] at (\p,0.96) {\lab}; }

%% ---------------- axe des pH ------------------------------------------------------
\draw[axe] (-0.6,-0.55) -- (14.9,-0.55) node[anchor=west, inner sep=3pt, font=\small] {pH};
\foreach \p in {0,1,...,14} { \draw[line width=0.4pt] (\p,-0.55) -- (\p,-0.44); }
\foreach \p in {0,7,14} {
  \draw[trait] (\p,0) -- (\p,-0.55);
  \draw[line width=0.7pt] (\p,-0.68) -- (\p,-0.44);
  \node[grad, anchor=north] at (\p,-0.70) {\p}; }
\node[petit, anchor=north, text=solideE, align=center] at (7,-1.02)
  {$\mathrm{pH} = 7{,}0$\\neutre};

%% ---------------- règle des concentrations ---------------------------------------
\draw[line width=0.6pt] (0,-1.75) -- (14,-1.75);
\foreach \p/\lab in {0/{$10^{0}$}, 7/{$10^{-7}$}, 14/{$10^{-14}$}} {
  \draw[line width=0.7pt] (\p,-1.87) -- (\p,-1.63);
  \node[grad, anchor=north] at (\p,-1.89) {\lab}; }
\node[petit, anchor=east] at (-0.5,-1.75) {$[\mathrm{H_3O^+}]$ (mol$\cdot$L$^{-1}$)};

%% ---------------- sens de variation de la concentration ---------------------------
\draw[axe, solideD, line width=0.9pt] (1,-2.85) -- (13,-2.85);
\node[petit, anchor=south, text=solideD] at (7,-2.82)
  {$[\mathrm{H_3O^+}] = 10^{-\mathrm{pH}}$ : la concentration \textbf{diminue} quand le pH augmente};
\end{tikzpicture}
\end{document}
